\begin{abstract}
We construct a singular normal affine complex surface whose tangent
sheaf is free of rank two. This disproves the Lipman--Zariski conjecture
in characteristic zero.
\end{abstract}

\section{Introduction}

For a finitely generated integral $\CC$-algebra $A$, write
\[
\Der_{\CC}(A)=\Hom_A(\Omega^1_{A/\CC},A).
\]
The module $\Der_{\CC}(A)$ consists of the algebraic vector fields on
$\Spec A$. On a smooth variety of dimension $d$, the tangent sheaf is
locally free of rank $d$. The Lipman--Zariski conjecture asks for the
converse: a complex algebraic variety with locally free tangent sheaf
should be nonsingular. In affine terms, it asserts that local freeness
of $\Der_{\CC}(A)$ forces $A$ to be regular. We prove the following.

\begin{theorem}\label{thm:main}
There exist a finitely generated normal integral $\CC$-algebra $A$ of Krull dimension two and a maximal ideal $\m\subset A$ such that
\[
\Der_{\CC}(A)\simeq A^{\oplus2},
\qquad A_{\m}\text{ is not regular}.
\]
\end{theorem}

Here freeness holds on the entire affine surface, which can be chosen
smooth away from the marked point. We first construct an isolated
Gorenstein analytic surface singularity. The passage to algebraic geometry
preserves its completed local ring, from which tangent freeness descends.

\subsection{Context and significance}

Lipman's study of free derivation modules established that local freeness
of the tangent sheaf forces normality~\cite[Theorem~3]{Lipman1965}.
The conjecture thus asks whether the same hypothesis excludes the remaining
normal singularities. In positive characteristic the corresponding assertion
fails even for hypersurfaces: over a perfect field of characteristic $p$,
the singular surface $xy-z^p=0$ has tangent sheaf freely generated by
$x\partial_x-y\partial_y$ and $\partial_z$
\cite[\S1]{GrafPositive2022}. The construction here has a different source
of tangent freeness and takes place entirely over $\CC$.

Several positive results explain both the strength of tangent freeness
and the role of surfaces in the problem. Hochster proved that a reduced
finitely generated characteristic-zero algebra with a nonnegative grading
and degree-zero part equal to the base field is a polynomial algebra if
its derivation module is free at the irrelevant maximal ideal
\cite[\S1, Theorem]{Hochster1977}. Van Straten and Steenbrink used
extension of differential forms to prove the isolated-singularity case
in dimension at least three~\cite[Theorem~1.3 and~\S1.6]{VanStratenSteenbrink1985}.
Their extension argument leaves the surface case as a distinct problem.

Hypersurfaces were treated by Scheja and Storch~\cite[Satz~9.3]{SchejaStorch1972};
K\"allstr\"om established the characteristic-zero algebraic locally complete-intersection
case~\cite[Corollary~1.3(1), corrected version]{Kallstrom2011}.
Extension theorems for differential forms also give strong results defined
by the singularities of the canonical divisor. Greb, Kebekus, Kov\'acs and
Peternell proved the conjecture for Kawamata log terminal spaces
\cite[Theorem~6.1]{GKKP2011}. The log-canonical case was proved by
Druel~\cite[Theorem~1.1]{Druel2014} and by Graf and Kov\'acs
\cite[Corollary~1.3]{GrafKovacs2014}.
These results relate tangent freeness to the behavior of differential
forms on a resolution. Another distinction concerns the Lie bracket:
J\"order proved smoothness for a normal complex space whose tangent sheaf
locally admits a basis of pairwise commuting vector fields
\cite[Theorem~1.1]{Jorder2014}. Thus no tangent frame near our singular
point can be chosen commuting.

For normal complex surfaces, the cohomology and configuration of the
exceptional curves give further criteria. Graf proved smoothness when
the geometric genus $p_g$ is at most one, and also when $p_g=2$ and the
exceptional divisor is not a tree of rational curves
\cite[Theorem~1.1]{GrafLowGenus2021}. Bergner and Graf then obtained the
criterion
\[
 p_g-g-b\le2,
\]
where $g$ is the sum of the genera of the exceptional curves and $b$ is
the first Betti number of the resolution's dual graph
\cite[Theorem~1]{BergnerGraf2020}.
The analytic singularity constructed here has one exceptional curve of genus $21$
and self-intersection $-10$. Its canonical discrepancy is $-5$, and
$p_g\ge33$, so $p_g-g-b\ge12$. Remark~\ref{analytic:invariants}
proves these calculations. They place the analytic example beyond the
log-canonical and geometric-genus criteria just described.

Global hypotheses provide other positive cases. Biswas, Gurjar and Kolte
prove smoothness for a normal affine surface whose tangent bundle on
the smooth locus is trivial and whose smooth locus has logarithmic
Kodaira dimension at most one~\cite[Theorem~0.1(1)]{BiswasGurjarKolte2014}.
Consequently, the smooth locus $U$ of the affine surface in
Theorem~\ref{thm:main} has logarithmic Kodaira dimension two
(Corollary~\ref{algebra:log-kodaira}). Here logarithmic Kodaira dimension
means the Iitaka dimension of $K_{\overline U}+D$ on a smooth projective
compactification with reduced simple-normal-crossing boundary $D$.
This global consequence uses the algebraic endpoint of the construction,
not just the analytic germ. In contrast, Graf's compact-surface result
identifies compact complex surfaces with free rank-two tangent sheaf as
complex two-tori~\cite[Corollary~1.3]{GrafLowGenus2021}.

The example also gives a negative answer in characteristic zero to the
derivation-blowup question of Barajas, Ch\'avez-Mart\'inez and
Romano-Vel\'azquez~\cite[Introduction, Question~1]{BarajasChavezRomano2026}.
The blowup at a module is the universal proper birational modification
making its pullback modulo torsion locally free. As those authors observe,
a singular variety with free tangent sheaf cannot be resolved by iterating
this operation at the module of derivations: every step is the identity.
For our normal surface, inserting normalization does not change this
conclusion (Corollary~\ref{algebra:derivation-blowups}).

The neighborhood of the exceptional curve is essential to the argument.
Hochster's theorem excludes obtaining the example from a singular
positively graded cone with free tangent sheaf. Instead, we start with a
split line bundle and change the gluing maps of its zero-section
neighborhood while preserving the normal bundle. This freedom supplies
the required vector fields. Their dual volume form has a fifth-order
pole along the exceptional curve, so holomorphic extension across that
curve is not an input to the construction.

\subsection{Main ideas and proof structure}

The construction first produces an analytic surface singularity and then
passes to an affine algebraic surface. The analytic step requires two
vector fields on a punctured neighborhood that remain linearly
independent everywhere. We build those fields on a resolution, where the
exceptional set is a smooth negative curve. Their determinant vanishes
along that curve to a prescribed order, and nowhere else near it.

Let $C$ be a smooth curve and $L$ a positive line bundle with
$L^{\otimes4}\simeq K_C$. Its dual bundle $L^{-1}$ provides the initial
neighborhood of $C$. In compatible base and fiber coordinates $(z,t)$,
that neighborhood carries the meromorphic two-form
\[
\sigma=t^{-5}\,dt\wedge dz.
\]
We seek two vector fields of the form
\[
v_i=t\bigl(A_i(t,z)t\partial_t+B_i(t,z)\partial_z\bigr),
\qquad i=1,2,
\]
with $\sigma(v_1,v_2)=1+O(t)$. The leading vertical coefficients form a
pair $r=(r_1,r_2)$ of sections of $L$ with six simple common zeros,
whose reduced divisor is denoted by $Z$. These common zeros create the
main compatibility problem: each requires a pointwise correction to make
the leading determinant coefficient nonzero along $C$, but the corrections must also satisfy
global residue constraints to glue the vector fields. The curve and trace
construction produces data satisfying both requirements.

\paragraph{A curve and normalized residues.}
Section~\ref{curve:section} constructs $C$ on $\PP^1\times\PP^1$, with
fourth-order contact at $Z$ with an auxiliary elliptic curve $\Gamma$.
A two-torsion relation on $\Gamma$ controls the section dimensions on $C$ and supplies
six vectors in a two-dimensional quadratic space. The resulting curve
has genus $21$, and $\deg L=10$.

Section~\ref{trace:section} transports these vectors to $C$ by a trace
pairing for a degree-four map $C\to\PP^1$. Residue duality compares
principal parts on the two curves and fixes the scalar and sign of the
pairing. Proposition~\ref{prop:residue-data} then gives a meromorphic
pair $P$ of sections of $L^2$ whose local expansions enforce the determinant
condition and whose residue sums cancel the gluing obstructions.

\paragraph{Finite vector-field jets.}
In Section~\ref{jets:section}, we modify the neighborhood by finitely
many holomorphic flows and choose compatible coefficients of the fields.
At order two the obstruction separates into a symmetric residue sum and
one antisymmetric sum; the normalized residue data cancel both.
Thereafter, we extend the vector fields and their desired determinant
as separate but compatible finite Taylor expansions. Choosing the
determinant coefficients before solving for the field coefficients leaves
the freedom needed to remove the remaining obstructions. Exactness of a
differential and cohomology vanishing finish the finite induction.
Proposition~\ref{prop:jets} supplies finite jets
with determinant $t^5(1+O(t))$ in ordinary coordinates, to sufficient
order for the holomorphic lifting below.

\paragraph{Holomorphic lifting and contraction.}
Section~\ref{analytic:section} realizes the finite transition maps as
a Hausdorff surface neighborhood $Y$ of $C$. Set
$\cI=\OO_Y(-C)$ and $\cF=\cT_Y(-\log C)$, the sheaf of vector fields
tangent to $C$. Grauert's criterion contracts the curve to a point
$x$ of a normal surface $X$~\cite{Grauert1962}. Cohomology of the
successive ideal quotients and the theorem on formal
functions~\cite{Grauert1960} give the cohomological vanishing needed to
lift these jets to holomorphic fields after shrinking, without requiring
an infinite sequence of neighborhood
changes to converge. Their determinant gives a frame on
$Y\setminus C\simeq X\setminus\{x\}$, which normality extends to a
free tangent sheaf on $X$. Adjunction proves that $x$ is singular.

\paragraph{The affine algebraic endpoint.}
Section~\ref{algebra:section} combines approximation with the finite
determinacy of equations together with their relations. The latter
method appears in Hironaka's equivalence theorem
\cite[\S3, Theorem~3.3]{Hironaka1969} and in Artin's algebraization
of isolated singularities~\cite[Theorems~1.10 and~3.8--3.9]{Artin1969}.
We prove the required finite-determinacy statement with a relation matrix:
the approximated equations and matrix retain their exact zero product.
This produces an algebraic local ring with the same completion as
$\OO_{X,x}$. A finite Jacobian presentation identifies its completed
derivation module with the continuous derivations of that completion.
Faithful flatness descends freeness and the local ring properties.
Finally, shrinking to one affine neighborhood gives the globally free
module in Theorem~\ref{thm:main}, while preserving the singular point.

\begin{table}[htbp]
\centering
\small
\begin{tabular}{@{}p{0.19\textwidth}p{0.73\textwidth}@{}}
\toprule
Notation & Meaning and normalization \\
\midrule
$C$, $Z$ & Smooth genus-$21$ curve and its six marked points. \\
$L$, $r$ & $\deg L=10$, $L^4\simeq K_C$; the two sections in $r$ have
common zero divisor $Z$. \\
$\Gamma$, $P$ & Auxiliary elliptic curve; meromorphic pair of sections
of $L^2$ supplying the residue identities. \\
$Y$, $\cI$, $\cF$ & Surface neighborhood of $C$; $\cI=\OO_Y(-C)$ and
$\cF=\cT_Y(-\log C)$. The jets lie in $\cF/\cI^9\cF$. \\
$\pi:Y\to X$ & Contraction of $C$ to the singular point $x$;
$C^2=-10$. \\
$A$, $\mathfrak p$ & Affine algebra and marked singular point;
\mbox{$\widehat{A_{\mathfrak p}}\simeq\widehat{\OO_{X,x}}$}. \\
\bottomrule
\end{tabular}
\caption{The main objects in the construction, in their order of use.}
\label{tab:notation}
\end{table}

\subsection{Conventions}

Line bundles on curves are written additively when convenient: $aL+bH$ means $L^{\otimes a}\otimes H^{\otimes b}$, and $H^q(L)$ means $H^q(C,L)$. A divisor and its associated line bundle may share a symbol. For two rows $X=(X_1,X_2)$ and $Y=(Y_1,Y_2)$, put
\[
X\wedge Y=X_1Y_2-X_2Y_1.
\]
The symmetric bilinear form associated with a quadratic form $Q$ is normalized by $\langle x,x\rangle=Q(x)$. A residue on an annular overlap means the contour integral divided by $2\pi i$. Logarithmic residue and adjunction use the convention that the normal differential comes first. All germs and varieties are over $\CC$. We use standard facts about algebraic curve cohomology and adjunction as in~\cite{Hartshorne1977}, and about coherent analytic sheaves and normal spaces as in~\cite{GrauertRemmert1984}.
